For \(r\in\{1,2\}\), let \(\Psi^{(r)}\in\mathbb S_+^p\) denote the invariant nuisance in the \(r\)-th fitted BACKSHIFT equation.  Thus, by the fitting hypotheses in the statement,
\[
D^{(r)}\Sigma_eD^{(r)\mathsf T}
 =\Psi^{(r)}+\operatorname{diag}(s^{(r)}_e),
 \qquad e\in H_r.
\]
First, inclusion--exclusion and \(|E|=m\), \(|H_r|\ge h=m-c\), give
\[
 |J|=|H_1\cap H_2|
 =|H_1|+|H_2|-|H_1\cup H_2|
 \ge 2h-m=m-2c=h-c.
\]

It remains to prove uniqueness.  Set \(Q=D^{(2)}(D^{(1)})^{-1}\).  The assumed noncollinearity implies that \(J\) is nonempty; fix \(e_0\in J\).  Subtracting the equation at \(e_0\) from the equation at \(e\in J\), and using the two displayed fits, yields
\[
 Q\operatorname{diag}(x_e)Q^{\mathsf T}
   =\operatorname{diag}(y_e),\qquad
 x_e=s^{(1)}_e-s^{(1)}_{e_0},\quad
 y_e=s^{(2)}_e-s^{(2)}_{e_0}.
\]
Let \(W=\operatorname{span}\{x_e:e\in J\}\subseteq\mathbb R^p\).  By linearity,
\[
 Q\operatorname{diag}(w)Q^{\mathsf T}\ \text{is diagonal for every }w\in W. \tag{1}
\]
For each \(k<\ell\), the coordinate projection of \(W\) onto coordinates \(k,\ell\) is all of \(\mathbb R^2\).  Indeed, if that projection had dimension at most one, some nonzero \((a,b)\in\mathbb R^2\) would satisfy
\[
 a(s^{(1)}_{ek}-s^{(1)}_{e_0k})
 +b(s^{(1)}_{e\ell}-s^{(1)}_{e_0\ell})=0
 \quad\text{for every }e\in J,
\]
placing the entire \(k,\ell\) cloud in one affine line, contrary to the hypothesis.

Put \(d_W=\dim(W)\).  Choose a basis \(w^1,\ldots,w^{d_W}\) of \(W\), and define its \(k\)-th coordinate profile by
\[
 u_k=(w^1_k,\ldots,w^{d_W}_k)^{\mathsf T}\in\mathbb R^{d_W}.
\]
The preceding projection assertion says that \(u_k,u_\ell\) are linearly independent whenever \(k\ne\ell\).  In particular, every \(u_k\) is nonzero.  Choose \(\beta\in\mathbb R^{d_W}\) outside the finitely many hyperplanes \(\{\beta:\beta^{\mathsf T}u_k=0\}\).  Having chosen \(\beta\), choose \(\alpha\in\mathbb R^{d_W}\) outside the finitely many hyperplanes
\[
 \left\{\alpha:\alpha^{\mathsf T}\bigl((\beta^{\mathsf T}u_\ell)u_k-(\beta^{\mathsf T}u_k)u_\ell\bigr)=0\right\},\qquad k<\ell.
\]
Every displayed normal vector is nonzero because \(u_k,u_\ell\) are linearly independent and both \(\beta^{\mathsf T}u_k\) and \(\beta^{\mathsf T}u_\ell\) are nonzero.  The choices exist because a finite union of proper linear hyperplanes cannot cover \(\mathbb R^{d_W}\): the product of their defining nonzero linear forms is a nonzero polynomial, while a nonzero real polynomial cannot vanish at every point (the latter follows by induction on the number of variables from the one-variable root bound).

Put
\[
 b=\sum_{q=1}^{d_W}\beta_qw^q,\qquad
 a=\sum_{q=1}^{d_W}\alpha_qw^q,\qquad
 \rho_k=\frac{a_k}{b_k}.
\]
The choice of \(\beta\) gives \(b_k\ne0\) for every \(k\), and the choice of \(\alpha\) gives \(\rho_k\ne\rho_\ell\) for \(k\ne\ell\).  By (1), both
\[
 G=Q\operatorname{diag}(a)Q^{\mathsf T}\quad\text{and}\quad
 K=Q\operatorname{diag}(b)Q^{\mathsf T}
\]
are diagonal.  Moreover, \(K\) is invertible because \(Q\) and \(\operatorname{diag}(b)\) are invertible.  Hence
\[
 GK^{-1}
 =Q\operatorname{diag}(\rho_1,\ldots,\rho_p)Q^{-1} \tag{2}
\]
is diagonal.  The eigenvalues on the right of (2) are pairwise distinct.  The eigenspaces of the diagonal matrix on the left are therefore the coordinate axes, so every column of \(Q\) is supported on exactly one coordinate.  Since \(Q\) is invertible, it is monomial: there are a permutation \(\pi\) and nonzero scalars \(q_i\) such that
\[
 Q_{i,\pi(i)}=q_i,\qquad Q_{ij}=0\quad(j\ne\pi(i)). \tag{3}
\]

It remains to remove the monomial ambiguity.  By \(\operatorname{diag}(D^{(1)})=\operatorname{diag}(D^{(2)})=\mathbf 1\), which is part of \(\text{def:admissible-matrices}\), equations \(D^{(2)}=QD^{(1)}\) and (3) imply
\[
 1=D^{(2)}_{ii}=q_iD^{(1)}_{\pi(i),i}. \tag{4}
\]
If \(\pi\) had a nontrivial cycle \(i_1,\ldots,i_d\), with \(d\ge2\) and \(\pi(i_t)=i_{t+1}\) cyclically, then (4) would give
\[
 \prod_{t=1}^d\left|D^{(1)}_{i_{t+1},i_t}\right|
 =\prod_{t=1}^d|q_{i_t}|^{-1}. \tag{5}
\]
The entries on the left of (5) form a directed simple cycle of \(I_p-D^{(1)}\), while
\[
 D^{(2)}_{i_t,i_{t+1}}
 =q_{i_t}D^{(1)}_{i_{t+1},i_{t+1}}=q_{i_t}
\]
form the reverse directed simple cycle of \(I_p-D^{(2)}\).  The two strict cycle-product inequalities in \(\text{def:admissible-matrices}\) would consequently imply both
\[
 \prod_{t=1}^d|q_{i_t}|^{-1}<1
 \quad\text{and}\quad
 \prod_{t=1}^d|q_{i_t}|<1,
\]
whose product is the contradiction \(1<1\).  Thus \(\pi\) is the identity.  Equation (4) and \(D^{(1)}_{ii}=1\) then give \(q_i=1\) for every \(i\).  Hence \(Q=I_p\) and \(D^{(1)}=D^{(2)}\), as claimed.